\documentclass[a4paper]{article}
\usepackage{amsfonts}
\usepackage{amsmath}
\usepackage{amssymb}
\usepackage{graphicx}

\begin{document}
  
\noindent \large Auteur: Siebe Mees \\
\noindent \large Studierichting: Informatica\\
\noindent \large Oefeningenreeks 7C nr. 1.5\\

\medskip

\normalsize

$f:\mathbb{R}^3 \rightarrow \mathbb{R}:(x,y,z) \mapsto \cos(xy)+\sin(yz)$\\

$\dfrac{\partial f}{\partial x}=-y\sin(xy)$\\

$\dfrac{\partial f}{\partial y}=-x\sin(xy)+z\cos(yz)$\\

$\dfrac{\partial f}{\partial z}=y\cos(yz)$\\

$\dfrac{\partial^2 f}{\partial x^2}=-y^2\cos(xy)$\\

$\dfrac{\partial^2 f}{\partial x \partial y}=-\sin(xy)-xy\cos(xy)$\\

$\dfrac{\partial^2 f}{\partial x \partial z}=0$\\

$\dfrac{\partial^2 f}{\partial y^2}=-x^2\cos(xy)-z^2\sin(yz)$\\

$\dfrac{\partial^2 f}{\partial y \partial x}=-\sin(xy)-xy\cos(xy)$\\

$\dfrac{\partial^2 f}{\partial y \partial z}=\cos(yz)-yz\sin(yz)$\\

$\dfrac{\partial^2 f}{\partial z^2}=-y^2\sin(yz)$\\

$\dfrac{\partial^2 f}{\partial z \partial x}=0$\\

$\dfrac{\partial^2 f}{\partial z \partial y}=\cos(yz)-yz\sin(yz)$\\

\begin{thebibliography}{999}
%\bibitem{a} Referentie
\end{thebibliography}
\end{document} 
